\section{Kiến thức nền tảng}
\label{sec:background}

Công nghệ container hóa đã thay đổi căn bản việc triển khai phần mềm đám mây hiện đại bằng cách cung cấp các môi trường thực thi nhẹ, có khả năng mở rộng và di động. Không giống như các máy ảo (VM), vốn đóng gói một hệ điều hành hoàn chỉnh cho mỗi phiên bản, các container chia sẻ nhân OS của máy chủ trong khi cung cấp các không gian người dùng cô lập thông qua không gian tên (namespaces) và nhóm điều khiển (cgroups). Kiến trúc này giảm chi phí vận hành và tăng tốc triển khai, khiến container trở nên lý tưởng cho microservices, quy trình CI/CD và khối lượng công việc linh hoạt. Tuy nhiên, mô hình chia sẻ nhân này mang lại những rủi ro bảo mật độc đáo---đặc biệt tại giao diện syscall, nơi điều phối quyền truy cập vào các chức năng cốt lõi của nhân.

\subsection{Bảo mật lời gọi hệ thống trong môi trường chia sẻ nhân}
\label{subsec:syscall_security}

Syscall đóng vai trò là giao diện chính giữa các ứng dụng không gian người dùng và tài nguyên nhân. Trong các hệ thống container hóa, chúng rất quan trọng đối với việc thực thi tiến trình, thao tác tệp, quản lý bộ nhớ và giao tiếp mạng. Thật không may, khả năng rộng lớn của chúng cũng khiến chúng trở thành mục tiêu hàng đầu của những kẻ tấn công. Nhân hệ điều hành cung cấp hàng trăm syscall, nhiều trong số đó là không cần thiết cho các khối lượng công việc container hóa điển hình, do đó mở rộng bề mặt tấn công~\cite{ref11_shu_security}.

Những kẻ tấn công thường khai thác các syscall bị lạm dụng hoặc có lỗ hổng để thực thi leo thang đặc quyền, thoát khỏi container hoặc can thiệp vào môi trường máy chủ. Các mối đe dọa như vậy trở nên nghiêm trọng hơn trong các triển khai đám mây đa thuê (multi-tenant), nơi nhiều container chia sẻ cùng một nhân máy chủ và do đó cùng một không gian tên syscall. Nếu không có kiểm soát syscall nghiêm ngặt, một sự xâm phạm trong một container có thể dẫn đến di chuyển ngang hoặc chiếm quyền máy chủ.

\subsection{Các loại mối đe dọa dựa trên lời gọi hệ thống}
\label{subsec:syscall_threats}

Lời gọi hệ thống mở ra các chức năng nhân cốt lõi, khiến chúng trở thành mục tiêu hàng đầu trong môi trường container hóa. Những kẻ tấn công khai thác chúng để leo thang đặc quyền, phá vỡ cô lập hoặc lạm dụng các hồ sơ quá rộng quyền.

\textbf{Leo thang đặc quyền:} Những kẻ tấn công có thể khai thác các lỗ hổng nhân, khả năng Linux bị cấu hình sai hoặc các hồ sơ seccomp không đầy đủ để đạt được đặc quyền cao hơn từ bên trong một container. Các ví dụ bao gồm việc sử dụng \texttt{setuid()}, \texttt{clone()} hoặc \texttt{execve()} theo các trình tự cụ thể để chiếm quyền thực thi luồng hoặc vượt qua các hạn chế không gian tên.

\textbf{Thoát container:} Các khai thác như \mbox{CVE-2022-0847}~\cite{cve_2022_0847} (Dirty Pipe) hoặc \mbox{CVE-2019-5736}~\cite{cve_2019_5736} (Docker runc) cho phép các tiến trình độc hại vượt qua ranh giới container và tương tác với máy chủ. Những điều này thường liên quan đến việc tạo ra các chuỗi syscall độc hại can thiệp vào bộ nhớ chỉ đọc hoặc thoát khỏi cô lập không gian tên.

\textbf{Lạm dụng bề mặt tấn công:} Nhiều container có quyền truy cập quá rộng vào các syscall do sử dụng hồ sơ seccomp chung hoặc mặc định. Sự quá rộng quyền này có thể bị các đối thủ tinh vi khai thác để kết hợp các syscall thành các cuộc tấn công nhiều giai đoạn và tàng hình.

\subsection{Hạn chế của các phương pháp giảm thiểu truyền thống}
\label{subsec:traditional_limitations}

Các cơ chế bảo mật như AppArmor, SELinux và Seccomp cung cấp lọc syscall và thực thi chính sách. Tuy nhiên, chúng thể hiện nhiều hạn chế trong môi trường đám mây hiện đại:

\begin{itemize}
  \item \textbf{Chính sách tĩnh không linh hoạt:} Các hồ sơ tĩnh không thích nghi với các hành vi mới hoặc
        vectơ mối đe dọa mới, thường dẫn đến kết quả dương tính giả (chặn quá mức) hoặc kết quả âm tính
        giả (bỏ sót tấn công).
  \item \textbf{Thiếu nhận thức về chuỗi:} Các hệ thống dựa trên quy tắc hoạt động trên các lời gọi
        syscall riêng lẻ mà không xem xét các phụ thuộc thời gian, khiến chúng mù quáng với các cuộc tấn
        công nhiều bước.
  \item \textbf{Nút thắt cổ chai về khả năng mở rộng:} Khi mật độ container tăng, các bộ lọc tĩnh trở nên
        khó duy trì và thực thi hiệu quả ở quy mô lớn.
  \item \textbf{Chi phí cấu hình thủ công:} Các chính sách phải được viết, kiểm tra và duy trì thủ công cho
        từng khối lượng công việc container---một gánh nặng không thực tế trong môi trường động.
\end{itemize}

Những điểm yếu này nhấn mạnh sự cần thiết của các giải pháp bảo mật syscall động, thông minh và nhận biết ngữ cảnh có thể mở rộng theo hạ tầng hiện đại.

\subsection{Tiềm năng của eBPF trong giám sát lời gọi hệ thống}
\label{subsec:ebpf}

eBPF là một công nghệ tích hợp nhân cho phép đo lường nhân an toàn, có thể lập trình cho các sự kiện nhân, bao gồm cả syscall~\cite{ref20_bpf_security}. Nó cho phép gắn kết động logic giám sát và lọc trực tiếp tại các điểm nhập syscall (ví dụ: \texttt{sys\_enter\_*}) với chi phí tối thiểu. Các lợi ích chính của eBPF trong lĩnh vực bảo mật syscall bao gồm:

\begin{itemize}
  \item \textbf{Thực thi thời gian thực:} Syscall có thể bị chặn hoặc kiểm tra ngay lập tức mà không cần
        chuyển đổi ngữ cảnh hay khởi động lại.
  \item \textbf{Lọc nhận biết ngữ cảnh:} Các chính sách có thể tính đến không gian tên, ID nhóm điều
        khiển, các đối số hoặc các mẫu syscall.
  \item \textbf{Khả năng mở rộng cao:} Các bản đồ và điểm theo dõi trong nhân hiệu quả cho phép thực thi
        trên hàng nghìn container đồng thời.
\end{itemize}

Tính chất có thể lập trình của eBPF khiến nó trở thành nền tảng lý tưởng để triển khai bảo mật syscall chi tiết, thời gian thực phù hợp với khối lượng công việc container.

\subsection{Phát hiện bất thường dựa trên chuỗi lời gọi hệ thống}
\label{subsec:sequence_anomaly}

Các mô hình học máy (ML), đặc biệt là những mô hình phân tích chuỗi syscall, cung cấp một bổ sung mạnh mẽ cho giám sát eBPF. Thay vì chỉ dựa vào các quy tắc được xác định trước, các mô hình dựa trên chuỗi có thể học các mẫu hành vi bình thường và gắn cờ các cửa sổ syscall bất thường lệch khỏi quy trình làm việc mong đợi.

Bộ mã hóa biến phân tự động (VAE)~\cite{vae_kingma} và Rừng cô lập~\cite{iforest_liu} (iForest) đặc biệt hiệu quả để mô hình hóa các chuỗi syscall:

\begin{itemize}
  \item VAE tái cấu trúc hành vi bình thường từ các biểu diễn không gian tiềm ẩn nén. Lỗi tái cấu trúc
        cao cho thấy sự bất thường.
  \item Rừng cô lập xác định các ngoại lệ thống kê dựa trên mức độ dễ dàng mà một chuỗi có thể được
        tách biệt khỏi các mẫu bình thường.
  \item Kết hợp cả hai phương pháp cải thiện tính mạnh mẽ bằng cách kết hợp tính điểm bất thường dựa
        trên tái cấu trúc và dựa trên cây.
\end{itemize}

Để hỗ trợ huấn luyện và xác nhận, các bộ dữ liệu công khai như bộ dữ liệu DongTing~\cite{ref26_dongtingdataset} cung cấp các chuỗi syscall được gán nhãn là bình thường hoặc độc hại. DeSFAM sử dụng độc quyền bộ dữ liệu DongTing, tận dụng 18.966 chuỗi được gán nhãn để huấn luyện và hiệu chỉnh các mô hình phát hiện bất thường của nó.

\subsection{Hướng tới bảo mật lời gọi hệ thống thích nghi và có thể mở rộng}
\label{subsec:adaptive_security}

Tính chất động của điều phối container đòi hỏi các cơ chế lọc syscall thích nghi với hành vi khối lượng công việc, tích hợp thông tin tình báo về mối đe dọa và mở rộng với sự can thiệp của con người ở mức tối thiểu.

DeSFAM giải quyết các nhu cầu này bằng cách kết hợp:
\begin{itemize}
  \item Lập hồ sơ syscall tĩnh-động lai.
  \item Phát hiện bất thường kết hợp học sâu (VAE) và các mô hình dựa trên cây (iForest).
  \item Thực thi chính sách dựa trên rủi ro sử dụng các móc eBPF LSM.
  \item Tính điểm rủi ro dựa trên khớp CVE và MITRE ATT\&CK.
  \item Các cơ chế tự tăng cường để tinh chỉnh hồ sơ seccomp tại thời điểm chạy.
\end{itemize}

Thiết kế tích hợp này cho phép phát hiện thời gian thực, chặn theo ngữ cảnh và giảm thiểu bề mặt syscall lâu dài, tạo tiền đề cho bảo mật container có khả năng phục hồi và mở rộng trong môi trường chia sẻ nhân.
